% !TEX root = ../main.tex
%% ---------------------------------------------------------------------------
%% Abstract and keywords
%% ---------------------------------------------------------------------------
\begin{abstract}
It is well known that numerical integration plays crucial roles in Galerkin meshfree methods for elasticity problems,
since inappropriate integration methods would lead to the loss of orthogonality condition, which is a basic part for evaluating the convergence performance.
As an extension of the works for second order problems, explicit estimates of $L_2$, $H_1$ and $H_2$ errors for Galerkin meshfree formulation to Kirchhoff plate (fourth order problem) was presented in this work.
Specifically, accounting for the effect of numerical integration, the $H_2$ error bound for orthogonally condition of Galerkin meshfree formulation was investigated, then it was regarded as a basic component to measure $H_2$ error of Galerkin meshfree methods.
Subsequently, with the help of closed graph theorem and Lagrange interpolation, the $L_2$ and $H_1$ error formulas of Galerkin meshfree methods was established and then divided into two components, which are respectively related to the standard error and integration error.
After that, the standard error was separately investigated with considering different orders of Lagrange interpolation for different orders of meshfree approximation.
It turns out that when the consistency of integral schemes satisfies the same order of meshfree approximation ($p=q$), an optimal convergence performance of meshfree methods to Kirchhoff plate will be achieved.
Furthermore, only when the order of meshfree approximation is greater than two ($p>2$) and the integral schemes possess the above property about consistency ($p=q$), a $(p+1)$th order of convergence rate can be realized.
Numerical examples well validate the proposed theory.
\end{abstract}

\begin{keywords}
Galerkin meshfree method \sep Kirchhoff plate \sep numerical integration \sep integration constraint \sep error estimate
\end{keywords}
